\documentclass[11pt]{article}
\usepackage[a4paper,margin=28mm]{geometry}
\usepackage{amsmath,amssymb,amsthm}
\usepackage[hidelinks]{hyperref}
\newtheorem{theorem}{Theorem}
\title{Fixed-power lower bounds for Hamiltonian Cycle, Hamiltonian Path, and Longest Path}
\author{Atlas\thanks{The proofs in this note were found using ChatGPT.}}
\date{18 November 2026}
\begin{document}
\maketitle

The Satisfiability Hypothesis (SH) states that, for every fixed
\(0<a<1\), no deterministic algorithm decides satisfiability of a
CNF formula of length \(s\) in time \(2^{\mathcal O(s^a)}s^{\mathcal O(1)}\)
~\cite{sh}. We use explicit inputs throughout. Counting bits or counting
literal occurrences, variables, and clauses gives the same family of
fixed-power exclusions after consecutively numbering the occurring
variables. Indeed the two lengths then differ by at most a logarithmic
factor, which a strictly larger exponent below one absorbs.
Let \(B\) denote the complete input length in bits.

The Hamiltonian Cycle exclusions follow from the stronger induced-grid
bound in Table~1 of \cite{sh}, which uses the construction of
de Berg et al.~\cite[Theorem~4.7]{grid}. For each forbidden power below
one, choose a sufficiently large fixed dimension. In fixed dimension
the number of edges is linear in the order. Splitting one vertex and
forcing the path endpoints gives the Hamiltonian Path exclusion, and
choosing the spanning-path budget gives the Longest Path exclusion.
We give an alternative proof using the reduction of \cite{ipz} and
prove the converse for Hamiltonian Cycle parameterized by graph size.

\hypertarget{hamiltonian-sh-forward}{}
\begin{theorem}
Assume SH. For every fixed \(0<a<1\), none of the following problems
admits a deterministic algorithm in time \(2^{\mathcal O(p^a)}B^{\mathcal O(1)}\).
\begin{itemize}
\item \textsc{Hamiltonian Cycle} in a simple undirected graph, with
\(p\) its number of vertices or its number of vertices plus edges.
\item \textsc{Hamiltonian Path} in a simple undirected graph, with \(p\)
its number of vertices.
\item \textsc{Longest Path}, with \(p\) the requested path length in
edges (equivalently, the requested number of vertices, up to one).
\end{itemize}
\end{theorem}
\begin{proof}
Split the long clauses of a CNF formula into clauses with at most three
literals. For a clause with literals \(\ell_1,\ldots,\ell_r\), \(r>3\),
replace it by
\[
 (\ell_1\lor\ell_2\lor y_1),\quad
 (\neg y_i\lor\ell_{i+2}\lor y_{i+1})\quad(1\le i<r-3),\quad
 (\neg y_{r-3}\lor\ell_{r-1}\lor\ell_r).
\]
The \(y_i\) are private to the clause. A satisfying assignment extends
by setting \(y_i\) true precisely when the first \(i+1\) literals are
all false. If every literal is false, the chain forces every \(y_i\)
true and violates the last clause. The resulting formula has
\(\mathcal O(s)\) clauses.

Impagliazzo, Paturi, and Zane~\cite[Section~3, following Theorem~5]{ipz}
give the polynomial-time reduction from 3-SAT with the number of clauses
as parameter to \textsc{Hamiltonian Circuit} with the number of edges as
parameter. The output has \(\mathcal O(s)\) edges. If using the directed
form, replace each vertex by a three-vertex path and attach incoming and
outgoing arcs to opposite ends. In a spanning undirected cycle the middle
vertex forces traversal of each path in full, and each joining edge links
an outgoing end to an incoming end. Traversing one path in its forward
direction fixes the forward direction on all paths. This gives the
undirected equivalence with linear overhead. In either form an isolated
vertex gives an immediate negative answer. Otherwise the number of
vertices is at most twice the number of edges. Thus both claimed
parameters for \textsc{Hamiltonian Cycle} are \(\mathcal O(s)\).
Bounded-size cases can be decided directly.

For \textsc{Hamiltonian Path}, split any vertex \(v\) of the undirected
cycle instance into two nonadjacent copies with the same neighbors and
attach a private leaf to each copy. A Hamiltonian cycle opens at \(v\)
to give a spanning path between the leaves. Conversely, the leaves must
be the two endpoints of any spanning path. Removing them and identifying
the copies closes the path to a Hamiltonian cycle of the original graph.
This applies for graphs with at least three vertices. Asking for a path
with one fewer edge than the number of vertices is exactly the same
decision in \textsc{Longest Path}.

Every construction therefore maps
the original formula to an equivalent instance of polynomial bit length
and with parameter \(\mathcal O(s)\). The proposed running time for any
one of the problems would contradict SH.
\end{proof}

\hypertarget{hamiltonian-size-sh-converse}{}
\begin{theorem}
The family of fixed-power exclusions for \textsc{Hamiltonian Cycle}
with \(p=|V|+|E|\) implies SH.
\end{theorem}
\begin{proof}
We encode a directed traversal of an undirected Hamiltonian cycle into
a CNF formula of length \(\mathcal O(p\log^2(p+2))\) bits. Give each
orientation of each edge a Boolean variable selecting that arc. Require
exactly one incoming and one outgoing selection per vertex. For a list
of \(d\) selections this takes \(\mathcal O(d+1)\) gates: scan the list,
maintaining whether some entry has been selected and rejecting a further
selection once the flag is set. Reject an isolated vertex directly.

Choose a root and give each vertex a binary position between zero and
\(|V|-1\), using \(\lceil\log_2 |V|\rceil\) bits. Require the root to
have position zero and every selected arc to advance the position by
one modulo \(|V|\). Comparing with \(|V|-1\) and performing a conditional
reset or binary increment takes \(\mathcal O(\log |V|)\) bounded-fan-in
gates per arc. Every selected directed cycle must have length divisible
by \(|V|\). The selected arcs cover the vertices by disjoint cycles, so
there must be one spanning cycle. For \(|V|\ge3\), it corresponds to an
undirected Hamiltonian cycle. Conversely, ordering a Hamiltonian cycle
from the root satisfies all constraints. Handle smaller inputs directly.

Encode each gate by a variable and its constant-size truth-table clauses
and require the output to be true. There are \(\mathcal O(p\log(p+2))\)
gates, so the stated bit bound includes the binary variable names. If SH
fails for some \(a<1\), use the resulting CNF algorithm and choose any
\(b\) with \(a<b<1\). Its running time on these formulas, including the
polynomial construction, is \(2^{\mathcal O(p^b)}B^{\mathcal O(1)}\),
contradicting the fixed-power family.
\end{proof}

\begin{thebibliography}{9}
\bibitem{ipz} R.~Impagliazzo, R.~Paturi, and F.~Zane.
Which problems have strongly exponential complexity?
\emph{Journal of Computer and System Sciences} 63 (2001), 512--530.
\href{https://doi.org/10.1006/jcss.2001.1774}{doi:10.1006/jcss.2001.1774}.
\bibitem{sh} MIT Hardness Group et al.
ETH is Dead, Long Live SH. 2026.
\href{https://arxiv.org/abs/2611.02650v2}{arXiv:2611.02650v2}.
\bibitem{grid}
M.~de Berg, H.~L.~Bodlaender, S.~Kisfaludi-Bak, D.~Marx, and T.~C.~van der Zanden.
\emph{A Framework for Exponential-Time-Hypothesis--Tight Algorithms and Lower Bounds in Geometric Intersection Graphs}.
SIAM Journal on Computing 49(6), 1291--1331, 2020.
\href{https://doi.org/10.1137/20M1320870}{doi:10.1137/20M1320870}.
\end{thebibliography}
\end{document}
